
\subsection{(\S6) The stretch of individual typologies in the formation of cultural attitudes}

\begin{tipbox}
This section aims to build the connection between cultural attitude
and psychological type, but I'm not quite interested in this, and hence only leaving
an opening paragrpah here.
\end{tipbox}

\begin{enumerate}

\item \ovalbox{\it Cultural attitudes}
%
In his book, Cultural Attitudes in Psychological Perspective, published after
two decades of work in 1984, Henderson identified four contrasting stances that
he saw as traditional orientations to culture, describing them as attitudes
that explain much about how different individuals choose to engage with the
culture in which they live. In this context, 'culture' can be thought of as the
expressions and products of the shared, rarely examined assumptions and values
of a group of people. These four stances that Henderson defined were the
social, the religious, the aesthetic, and the philosophic.

\end{enumerate}
